% Per-game replacements of the fixed texts of lang-en.tex (\renewcommand).
\renewcommand{\tComputerText}{\textbf{Any MSX1} with two cartridge slots, one for the game and one for the V9990, and a third one if you add a Moonsound. \tExpanderText}
\renewcommand{\tSoundText}{The MSX PSG is used. A \textbf{Moonsound} (OPL4) is optional: it adds arcade music and voices. With more than 128~KB of Moonsound sample RAM, some effects also use arcade sounds.}
\renewcommand{\tControlsText}{The keyboard, for one or two players. Joysticks are optional: port 1 for player 1, port 2 for player 2.}
\renewcommand{\tSetV}{Insert the \textbf{V9990 cartridge} in a free slot and, if you have one, the \textbf{Moonsound} in another. \tExpanderText}
\renewcommand{\tSetGame}{Insert the \textbf{\gameTitle{} cartridge} in another free slot. The order does not matter.}
\renewcommand{\tSetJoy}{Plug joysticks into port 1 (player 1) and port 2 (player 2) if you want to use them; the keyboard works for both players anyway.}
\renewcommand{\tSetOn}{Switch the MSX on. The DiHalt Studio logo appears on the V9990's monitor. If the MSX's own screen says \texttt{This cartridge needs a V9990.} the V9990 is missing or not seated, and the game stops there.}
\renewcommand{\tPowerOnCaption}{Left: the cover picture. Right: the title screen.}
\renewcommand{\tCreditsTitle}{11. Credits and rights}
\newcommand{\tCreditsExtra}{\item \textbf{\gameTitle} was first an arcade game by \textbf{Atari Games}, 1987. The title screen of the MSX version credits \texttt{Z80 BY B. COSTAS} and \texttt{68000 BY J.TRIPP}.
\item With a Moonsound, the music and voices are those of the arcade game, adapted; the song of the logo screen is by DiHalt Studio.}
